% sec-07: UC San Diego Moores Cancer Center.  No figure; this section is
% deliberately prose and table only, because the institutional route is already
% carried by application 10's figures and repeating them here would duplicate.
\section{UC San Diego Moores Cancer Center}

\draftinstr{full-apply writes 4,500 to 5,500 characters from
\appfile{funding/potential-partners/UC-San-Diego/README.md} and
\appfile{funding/potential-partners/UC-San-Diego/priority-steps.md}, plus
\appfile{funding/potential-partners/Scripps/} for the alternate-route contrast.
No figure is placed in this section: application 10 already carries the intake
sequence, the authority diagram, the feasibility grid, and the approval
dependency, and repeating any of them here would be duplication rather than a
new perspective.}

\subsection{The Positioning Constraint}

\draftinstr{One paragraph stating the constraint in full: UC San Diego is named
as the intended partner of choice and nothing more; it is not described as a
partner, sponsor, trial site, or endorser, and will not be without written
authorization.}

\subsection{Three Corrections the Site Research Forced}

\draftinstr{full-apply builds Table 8: columns Claim as first drafted,
Correction, and Where it now appears. Three rows, from
\appfile{priority-steps.md} \S1: daraxonrasib is not first-in-human; the robotic
configuration must be named by manufacturer, model, cleared configuration, arm
count, instruments and software version; no drug supply or letter of
authorization exists. State that nine of the ten applications carried the first
error until the PART I audit pass.}

\subsection{Two Routes, Sequenced Differently}

\draftinstr{full-apply builds Table 9 comparing the UC San Diego route with the
Scripps route on four rows: entry point, first ask, the step at which the trial
appears, and the risk if the route fails. Source both columns from the two site
READMEs.}
